\documentclass[10pt,a4paper]{amsart}
\usepackage[margin=19mm]{geometry}
\usepackage{amsmath,amssymb,mathtools}
\usepackage[scr=rsfs,cal=euler]{mathalfa}
\usepackage{xfrac}
\usepackage[unicode,hidelinks,bookmarks=false]{hyperref}
\newcommand{\ie}{i.e.\ }
\newcommand{\cf}{cf.\ }
\newcommand{\resp}{respectively\ }
\newcommand{\Subsection}{\subsection}
\providecommand{\nobreakdash}{\nobreak\hbox{-}}
\setlength{\emergencystretch}{2em}
\multlinegap0pt
\usepackage{mathtools}
\usepackage{amssymb}
\usepackage{tcolorbox}
\usepackage{xcolor}
\usepackage{tikz}
\usepackage{enumerate}
\newtheorem{thm}{Theorem}
\newtheorem{lem}{Lemma}
\title{Lie algebroid Connections, Moduli of $\mathcal{L}$--twisted Principal Objects and motives}
\author{Samit Ghosh, Arjun Paul}
\date{Source: arXiv:2605.18304v1 (CC BY 4.0)}
\begin{document}
\maketitle

\noindent\emph{Source and licence.} Lie algebroid Connections, Moduli of $\mathcal{L}$--twisted Principal Objects and motives. Version arXiv:2605.18304v1 (CC BY 4.0), distributed by arXiv under the Creative Commons Attribution 4.0 International licence, http://creativecommons.org/licenses/by/4.0/, as recorded in the arXiv licence metadata for this exact version and shown on the abstract page https://arxiv.org/abs/2605.18304v1. Full licence notice: \texttt{licenses/arxiv-2605.18304-CC-BY-4.0.txt}.\par\smallskip

\noindent\textit{Editorial scope.} Counted unit: the article's thm L-connection (source0.tex lines 2135--2145). The article's complete original proof of it is delivered with it (source0.tex lines 2147--2238, one \texttt{proof} environment of 92 lines). No mathematics is written, repaired or re-derived: the statement and the proof are the article's own lines, in the article's own notation, and no retained line is substituted or re-flowed.  Retained uncounted as the source-local context the proof needs: lines 897--913; lines 1793--1811; lines 1902--1925; lines 2001--2018.  Nothing was omitted from any retained span, so the package carries no omission record.  No retained line contains a citation, so the wrapper carries no bibliography and no bibliography entry.  No retained line contains a figure environment, an image-inclusion command or drawing code, so no image is delivered and the package declares no asset dependency.  Editorial formatting only, in addition: an AMS \texttt{amsart} preamble that restates the article's own declarations and macros used by the retained lines, namely the environments \texttt{thm}, \texttt{lem} on separate counters, together with the standard packages those lines need.  The retained environments print in this document as Theorem 1, Lemma 1 in order of appearance, whereas the article prints the same items with the numbering of its own full document.  The complete native source of the article as distributed in the arXiv e-print for this exact version is bundled unchanged as \texttt{sources/200847/source0.tex}.
\begin{thm}\label{eqn: moduli-of-L-connection on VB}
Let $\mathcal{L}$ be a transitive Lie algebroid and let $X$ be a smooth projective variety over $\mathbb{C}$. There exists a
quasi-projective scheme $\mathcal{M}^{\text{DR}}_{\mathcal{L}}(X,n)$ over $\mathbb{C}$ which universally corepresents the functor 
$\underline{\mathcal{M}^{\text{DR}}_{\mathcal{L}}}(X,n),$
assigning to a $\mathbb{C}$--scheme $S$ the set of isomorphism classes of
vector bundles $\mathcal{E}$ of rank $n$ on $X' := X \times S$
equipped with integrable $\mathcal{L}$--connections, and having Hilbert
polynomial $n\mathcal{P}_{0}$.
	
Fix a base point $\mathfrak{p} \in X$.  
Then there exists a quasi-projective scheme $\mathcal{R}^{\text{DR}}_{\mathcal{L}}(X,\mathfrak{p},n)$ over $\mathbb{C}$ representing the functor which assigns to a $\mathbb{C}$--scheme $S$ the set of isomorphism classes of triples $(\mathcal{E},D_{\mathcal{L}},\beta)$, where
\begin{itemize}
		\item $(\mathcal{E},D_{\mathcal{L}})$ is a vector bundle of rank $n$ on $X' = X \times S$ equipped with an integrable $\mathcal{L}$--connection, and
		\item $\beta : \mathcal{E}|_{\mathfrak{p}} \xrightarrow{\sim} \mathbb{C}^{n}$ is a framing at the point $\mathfrak{p}$.
\end{itemize}
	Moreover, with respect to a suitable linearization, all points of $\mathcal{R}^{\text{DR}}_{\mathcal{L}}(X,\mathfrak{p},n)$ are semistable for the natural action of $GL_{n}(\mathbb{C})$, and the resulting good quotient is naturally identified with $\mathcal{M}^{\text{DR}}_{\mathcal{L}}(X,n)$.
\end{thm}

\begin{lem}\label{lem:projective-scheme-N(E,k)}
	Let $\mathcal{L}$ be a transitive Lie algebroid on $X$, and let $(E,\nabla_{\mathcal L})$
	be a vector bundle equipped with an integrable $\mathcal L$--connection.
	Fix an integer $k$.
	Then there exists a projective scheme $\mathcal{N}(E,k)$ over $\mathbb{C}$
	representing the functor which assigns to each $\mathbb{C}$--scheme $S$
	the set of quotients
	\[
	p_X^*E \twoheadrightarrow F \to 0
	\quad \text{on } X_S := X\times S,
	\]
	which are compatible with the $\mathcal L$--connection.
	Moreover, the natural morphism
	\[
	\mathcal{N}(E,k) \longrightarrow
	\mathfrak{Grass}\bigl(E|_{\mathfrak p\times S},k\bigr)
	\]
	is a closed embedding.
\end{lem}

\begin{lem}\label{monodromy related thm-for-PB with-L-connection}
	Let $G \subset H$ be algebraic groups. Let $X$ be a smooth projective variety
	and let $\mathcal{L}$ be a transitive Lie algebroid on $X$.
	Suppose $E_{G}$ is a principal $G$--bundle on $X$ equipped with an integrable
	$\mathcal{L}$--connection. Then the associated bundle
	\[
	E_{H} := E_{G} \times^{G} H
	\]
	admits a natural structure of a principal $H$--bundle with integrable
	$\mathcal{L}$--connection.
	
	Moreover, fixing a closed point $\mathfrak{p} \in X$, this construction induces
	a bijection between the following sets:
	\begin{enumerate}
		\item isomorphism classes of pairs $(E_{G},\mathfrak{b}')$, where $E_{G}$ is a
		principal $G$--bundle with integrable $\mathcal{L}$--connection and
		$\mathfrak{b}' \in (E_{G})_{\mathfrak{p}}$ is a point;
		
		\item isomorphism classes of pairs $(E_{H},\mathfrak{b})$, where $E_{H}$ is a
		principal $H$--bundle with integrable $\mathcal{L}$--connection and
		$\mathfrak{b} \in (E_{H})_{\mathfrak{p}}$ is a point, such that the monodromy of
		$(E_{H},\mathfrak{b})$ is contained in $G$.
	\end{enumerate}
\end{lem}

\begin{lem}\label{lem:Fr-bundle-induce-L-connection}
	Let $E$ be a vector bundle of rank $n$ on $X$ equipped with an integrable
	$\mathcal{L}$--connection, and let
	\[
	\beta \colon E|_{\mathfrak{p}} \xrightarrow{\;\sim\;} \mathbb{C}^{n}
	\]
	be a choice of framing at a closed point $\mathfrak{p} \in X$.
	Then the frame bundle $\mathbf{Fr}(E)$ carries a natural structure of a
	principal $GL(n,\mathbb{C})$--bundle with integrable $\mathcal{L}$--connection.
	Moreover, the vector bundle $E$ is recovered as the associated bundle
	\[
	E \;=\; \mathbf{Fr}(E) \times^{GL(n,\mathbb{C})} \mathbb{C}^{n}.
	\]
	This construction induces a bijection between the sets of isomorphism classes
	of pairs $(E,\beta)$ and $(\mathbf{Fr}(E),\mathfrak{b})$, where
	$\mathfrak{b} \in \mathbf{Fr}(E)|_{\mathfrak{p}}$ corresponds to the
	framing $\beta$.
\end{lem}

\begin{thm}\label{thm: L-twisted - reprsentaion - space for L-connection}
	Let $\mathfrak{p}$ be a fixed base point of $X$. Then there exists a scheme $\mathcal{R}^{\text{DR}}_{\mathcal{L}}(X,\mathfrak{p},G)$
	representing the functor
	$\mathcal{M}^{\text{DR}^{\natural}}_{\mathcal{L}}(X,\mathfrak{p},G)$.
	Moreover, if $f \colon G \hookrightarrow H$ is a closed embedding of algebraic groups, then $f$ induces a closed embedding
	\[
	\mathcal{R}^{\text{DR}}_{\mathcal{L}}(X,\mathfrak{p},G)
	\hookrightarrow
	\mathcal{R}^{\text{DR}}_{\mathcal{L}}(X,\mathfrak{p},H).
	\]
\end{thm}

\begin{proof}
	We first treat the case $G = GL(n,\mathbb{C})$.
	By Lemma~\eqref{lem:Fr-bundle-induce-L-connection}, framed vector bundles with integrable $\mathcal{L}$--connections are equivalent to  framed principal $GL(n,\mathbb{C})$--bundles with integrable $\mathcal{L}$--connections. Hence, by Theorem~\eqref{eqn: moduli-of-L-connection on VB}, the functor
	$\mathcal{M}^{\text{DR}^{\natural}}_{\mathcal{L}}(X,\mathfrak{p},GL(n,\mathbb{C}))$ is represented by the scheme
	\[
	\mathcal{R}^{\text{DR}}_{\mathcal{L}}(X,\mathfrak{p},GL(n,\mathbb{C}))
	:= \mathcal{R}^{\text{DR}}_{\mathcal{L}}(X,\mathfrak{p},n).
	\]
	
	Now let $G \subset H$ be a closed algebraic subgroup, and assume that the scheme $\mathcal{R}^{\text{DR}}_{\mathcal{L}}(X,\mathfrak{p},H)$ representing
	$\mathcal{M}^{\text{DR}^{\natural}}_{\mathcal{L}}(X,\mathfrak{p},H)$ is already constructed.
	Let $
	(\mathcal{F}^{\mathrm{univ}}, \nabla^{\text{DR}^{\mathrm{univ}}}_{\mathcal{L}},b^{\mathrm{univ}})
	$
	denote the universal framed principal $H$--bundle with integrable
	$\mathcal{L}$--connection on
	$X \times \mathcal{R}^{\text{DR}}_{\mathcal{L}}(X,\mathfrak{p},H)$.
	
	Let $V$ be a finite--dimensional representation of $H$, and let
	$W \subset V$ be a subspace preserved by $G$.
	Set
	\[
	E^{\mathrm{univ}} := \mathcal{F}^{\mathrm{univ}} \times^{H} V,
	\]
	which is a vector bundle with integrable $\mathcal{L}$--connection.
	The framing $b^{\mathrm{univ}}$ induces an isomorphism
	\[
	\beta^{\mathrm{univ}} :
	E^{\mathrm{univ}}|_{\mathfrak{p} \times \mathcal{R}^{\text{DR}}_{\mathcal{L}}(H)}
	\;\xrightarrow{\;\sim\;}\;
	\mathcal{V} := V \otimes_{\mathbb{C}}
	\mathcal{O}_{\mathcal{R}^{\text{DR}}_{\mathcal{L}}(H)}.
	\]
	
	Let $k := \dim(V) - \dim(W)$, and denote by
	$\mathcal{W} := W \otimes_{\mathbb{C}}
	\mathcal{O}_{\mathcal{R}_{\mathcal{L}}(H)}$
	the corresponding trivial subbundle of $\mathcal{V}$.
	The quotient $\mathcal{V} \to \mathcal{V}/\mathcal{W}$ defines a section
	\[
	\sigma_{\mathcal{V}/\mathcal{W}} :
	\mathcal{R}^{\text{DR}}_{\mathcal{L}}(H)
	\longrightarrow
	\mathfrak{Grass}_{\mathcal{R}^{\text{DR}}_{\mathcal{L}}(H)}(\mathcal{V},k).
	\]
	
	By Lemma~\eqref{lem:projective-scheme-N(E,k)}, there exists a
	closed subscheme
	\[
	\mathcal{N}(E^{\mathrm{univ}},k)
	\subset
	\mathfrak{Grass}_{\mathcal{R}^{\text{DR}}_{\mathcal{L}}(H)}
	\bigl(E^{\mathrm{univ}}|_{\mathfrak{p}\times
		\mathcal{R}^{\text{DR}}_{\mathcal{L}}(H)}, k\bigr),
	\]
	parametrizing quotients compatible with the integrable
	$\mathcal{L}$--connection.
	Using the framing $\beta^{\mathrm{univ}}$, we view
	$\mathcal{N}(E^{\mathrm{univ}},k)$ as a closed subscheme of
	$\mathfrak{Grass}_{\mathcal{R}^{\text{DR}}_{\mathcal{L}}(H)}(\mathcal{V},k)$.
	Define
	\[
	\mathbf{C}(V,W)
	:=
	\sigma_{\mathcal{V}/\mathcal{W}}^{-1}
	\bigl(\mathcal{N}(E^{\mathrm{univ}},k)\bigr)
	\;\subset\;
	\mathcal{R}^{\text{DR}}_{\mathcal{L}}(H).
	\]
	
	By construction, a morphism $g \colon S \to \mathcal{R}^{\text{DR}}_{\mathcal{L}}(H)$ factors
	through $\mathbf{C}(V,W)$ if and only if the pullback
	$g^{*}(E^{\mathrm{univ}})$ admits a strict subbundle preserved by the
	$\mathcal{L}$--connection whose fiber at $\mathfrak{p}$ corresponds to $W$.
	Equivalently, this condition expresses that the monodromy of the framed object
	$g^{*}(\mathcal{F}^{\mathrm{univ}}, b^{\mathrm{univ}})$ preserves $W$.
	
	Finally, define
	\[
	\mathcal{R}^{\text{DR}}_{\mathcal{L}}(X,\mathfrak{p},G)
	:=
	\bigcap_{(V,W)} \mathbf{C}(V,W),
	\]
	where the intersection is taken over all representations $V$ of $H$ and all $G$--invariant subspaces $W \subset V$.
	This is a closed subscheme of
	$\mathcal{R}^{\text{DR}}_{\mathcal{L}}(X,\mathfrak{p},H)$.
	By construction, it represents the functor assigning to a scheme $S$ the set of framed principal $H$--bundles with integrable $\mathcal{L}$--connection whose monodromy is contained in $G$.
	By Lemma~\eqref{monodromy related thm-for-PB with-L-connection}, this functor identifies with $\mathcal{M}^{\text{DR}^{\natural}}_{\mathcal{L}}(X,\mathfrak{p},G)$.
	
	The final assertion concerning closed embeddings follows immediately from the
	construction.
\end{proof}

\end{document}
